\refstepcounter{section}
\section*{Lecture 16: Introducing the Dirac Equation}
\addcontentsline{toc}{section}{Lecture 16: Introducing the Dirac Equation}
We saw that, when quantised, the Klein-Gordon equation provided a negative probability density, which was problematic since it had no obvious interpretation. Hence we said that, the KG equation cannot describe a single-particle system, instead,it describes a spin-0 particle \textit{field}. All the problems arose since, from the relativistic dispersion relation, \begin{equation}\label{reldisp}
    E^2=p^2+m^2 = p_x^2+p_y^2+p_z^2+m^2
\end{equation} we obtain a differential equation which is second order in both space and time. \\[0.2cm]
Let us now look for something which we force to be first order in space and time from the beginning. The most sane choice is to just take the square-root of the dispersion relation, however, square-root of an operator is ill-defined and has many other subtleties. We want everything to be as simple as possible. Hence, we try to \textit{linearise} the square-root. We thus write,
$$E = \sqrt{p^2+m^2} = \alpha_1 p_1 + \alpha_2 p_2+\alpha_3 p_3 + \beta m$$
We first promote the momentum and energy to operators, that is, $p\to -i \nabla$ and $E\to i \partial_t $. Then the Schr\"odinger equation becomes, 
$$i \partial_t \psi = \qty(-i\alpha_i  \pdv{x^i} + \beta m)\psi \qquad \iff \qquad H\psi = E\psi$$ Since we want this to match with the actual dispersion relation Eq. \ref{reldisp}, let us see what happens when we apply the Hamiltonian twice on the time-independent equation  (which would give us an $E^2$ ).\\[0.2cm]
Note that, since the \textit{linearisation} of the square-root is extremely non-trivial, $\alpha_i$ and $\beta$ might not be just numbers, they could be literally anything. Thus, we do not make any apriori assumptions!\footnote{Well, we do assume that $\alpha_i$ and $\pdv{x^i}$ commute in some sense}
\begin{align*}
    H^2 \equiv &= \qty(-\alpha_i^2 \pdv[2]{{x^i}}+\beta m^2 - (\alpha_i \alpha_j + \alpha_j \alpha_i)\frac{\partial^2}{\partial {x^i} \partial {x^j}} - im (\beta a_i + \alpha_i \beta)\pdv{x^i})
\end{align*}
If the relativistic dispersion relation has to be satisfied, we find some constraints on $\alpha_i$ and $\beta$:
\begin{alignat*}{3}
\alpha_i^2 &= \mathds{1}\qquad  \alpha_i \alpha_j + \alpha_j \alpha_i &&= 2\delta_{ij}\mathds{1}\\
\beta^2 &=\mathds{1} \qquad \beta \alpha_i + \alpha_i \beta &&= 0
\end{alignat*}
where $\mathds{1}$ represents the identity operator in appropriate space. Having $\alpha_i$ and $\beta$ to be just simple scalar numbers clearly do not satisfy the equations, hence let us assume for now, that these are matrices. We then additionally impose the contraint that $\alpha_i^\dagger = \alpha_i$ and $\beta^\dagger = \beta$ since energy and momentum operator (and hence the Hamiltonian) are Hermitian. . With this assumption, let us note some nice properties of these matrices. 
\begin{itemize}
    \item Let $v$ be an eigenvector of $\alpha_i$ $\implies \alpha_i v = \lambda v\implies \underbrace{\alpha_i^2}_{\mathds{1}} v = \lambda^2 v \implies \lambda = \pm 1$ . Similar argument goes for $\beta$. Thus, these matrices have eigenvalues either +1 or -1. 
    \item $\alpha_i\beta = -\beta \alpha_i \implies \alpha_i\beta^2 = -\beta\alpha_i\beta \implies \alpha_i = \beta\alpha_i\beta\implies \Tr \alpha_i = -\Tr \alpha_i \implies \Tr \alpha_i=0 $,
    where we have used the cyclic property of trace. Similar argument goes for $\beta$. Hence we obtain that these matrices are also traceless.
    \item Since thre trace, that is, the sum of the eigenvalues, should equal zero, we should definitely have half of the eigenvalues as +1 and the other half to be -1. Thus, these matrices are guaranteed to be \textit{even} dimensional. 
    \item The relations between $\alpha_i$  looks suspiciously similar to the Pauli matrices and we can choose $\alpha_i$ to be $\sigma_i$. Alas, we run into a dangerous problem. We note that $\beta$ cannot be the a two-dimensional matrix, since then it would not be consistent with the anticommutation relation with $\alpha_i\equiv \sigma_i$. Thus, the next option is to try for \textit{four dimensional} matrices.
\end{itemize}
\pagebreak
Let us define the following, 

$$\alpha_i := 
\begin{bNiceArray}{c:c}[margin,columns-width=auto]
  0 & \Block{1-1}{\sigma_i}\\
\hdashline
  \sigma_i & \Block{1-1}{0}\\
\end{bNiceArray} \qquad \beta := 
\begin{bNiceArray}{c:c}[margin,columns-width=auto]
  \mathds{1} & \Block{1-1}{0}\\
\hdashline
  0 & \Block{1-1}{-\mathds{1}}\\
\end{bNiceArray} $$
With this definition, we can indeed verify that all the conditions hold \footnote{Note that these matrices are $4\times 4$ matrices. Each element in the matrix are indeed $2\times 2$ matrices. } and we obtain a way to proceed further. We have an expression for the equation as,
$$i \pdv{\psi}{t} = -i \alpha\cdot \nabla + \beta m \implies i \beta\pdv{\psi}{t} = -i\beta \alpha\cdot \nabla + \underbrace{\beta^2}_{\mathds{1}} m$$
Let us define the following matrices:
\begin{align*}
  \gamma^0 &:= \beta\\
\gamma^i &:=\beta\alpha_i
\end{align*}
Then, using this definition in the above expression, we have:
\begin{align*}
  i \gamma^0 \partial_0\psi  +i \gamma^i \partial_i\psi - m\psi =0 \implies (i \gamma^\mu \partial_\mu - m)\psi =0
\end{align*}
Note that, we had introduced the 4-vector concept in an ad-hoc manner. However, doing this, we finally derive the well-known form of the \textit{Dirac equation}!
